\originalsection{18.433 Fall 2003：课程作业}
本节收录 Santosh Vempala 课程的图论相关作业：作业1全部七题、作业2全部五题、作业3第6题、作业4第2、3题。保留原大题和字母小问，共15大题、25个叶题。官方公开集合没有独立答案文件，以下均为编者独立解答。通常图论题采用有限简单图；流网络的每条原弧保留独立身份。个别原题依赖课堂约定，所需补条件逐题注明。整数规划、极点和分离预言机在相关解答中交代；精确多项式线性优化算法的基础定理另列调用边界。

\par\addvspace{7pt}\Needspace{4\baselineskip}\noindent\textbf{作业1：匹配理论}

\begin{exercise}\label{mit:18433-a1-q1}
\textbf{18.433 Fall 2003，作业1，题1.}
证明：一个图是二分图，当且仅当它不含奇数长度的圈。
\end{exercise}
\sourceof{\href{https://ocw.mit.edu/courses/18-433-combinatorial-optimization-fall-2003/resources/a1/}{官方作业1}，PDF第1页；编者独立解答，非官方答案}
\begin{solution}
若图的顶点分为 $A,B$，每条边跨两侧，则沿圈每走一步换一侧，回到起点必须走偶数步，所以没有奇圈。

反过来，在每个连通分量中选根 $r$，按到 $r$ 的最短路长度的奇偶给顶点分两侧。若某条边 $uv$ 的两端距离同奇偶，取两条根到端点的最短路，再接边 $uv$，得到奇长闭行走。说明奇长闭行走一定含奇圈：若行走中除首尾外有重复顶点，就在该顶点处分成两段闭行走；总长度为奇数，至少一段仍为奇数。反复取这一较短奇段，最后得到没有内部重复顶点的奇圈。它与假设矛盾。因此每条边两端距离异奇偶，各分量的二分合并便给出全图二分。
\end{solution}

\begin{exercise}\label{mit:18433-a1-q2}
\textbf{18.433 Fall 2003，作业1，题2.}
考虑寻找最大匹配的贪心过程：先任取一条边作为匹配，不断加入与当前匹配没有公共端点的边，直到不能再加。
\begin{enumerate}[label=(\alph*)]
\item 对有 $n$ 个顶点、$m$ 条边的图，分析运行时间。
\item 给出它没有找到最大匹配的图。
\item 证明所得匹配的大小总至少是最大匹配大小的一半。
\end{enumerate}
编者补充：无边图返回空匹配；第(a)问的时间界按下述一次扫描的实现计算。
\end{exercise}
\sourceof{\href{https://ocw.mit.edu/courses/18-433-combinatorial-optimization-fall-2003/resources/a1/}{官方作业1}，PDF第1页；编者独立解答，非官方答案}
\begin{solution}
(a) 为每点维护一个是否已被匹配的标记，初始为否。按任意顺序扫描全部边，遇到两端均未标记的边就选入并标记两端。每条边只查两次标记，总时间 $O(n+m)$，额外空间 $O(n)$，输出最多 $\lfloor n/2\rfloor$ 条边。此前跳过的边至少已有一端被标记，以后也不会变为可加入，因此一次扫描所得确为极大匹配。若每轮从头搜索可加入边，时间可以较大；贪心规则本身并未指定唯一实现。

(b) 四顶点路 $1-2-3-4$ 若先选 $23$，便不能再选其他边，输出大小1；最大匹配 $\{12,34\}$ 大小2。

(c) 设输出为 $M$，任取最大匹配 $M^*$。极大性使 $M$ 的全部端点覆盖每一条边，否则两个自由端点之间的边还可以加入。故每条 $M^*$ 边至少占用 $M$ 的一个端点；$M^*$ 的边彼此不共端点，不同边占用的端点不同。可占用端点共 $2|M|$ 个，所以 $|M^*|\le2|M|$。第(b)问还表明这一保证不能普遍提高。
\end{solution}

\begin{exercise}\label{mit:18433-a1-q3}
\textbf{18.433 Fall 2003，作业1，题3.}
在边加权图中，每次贪心加入与当前匹配不相交的最重边，直到不能再加。证明所得匹配总权至少为最优匹配总权的一半。

编者补条件：边权非负。原题没有把这一课堂约定写在题面中；若允许负权且要求一直加到极大，原结论不成立。
\end{exercise}
\sourceof{\href{https://ocw.mit.edu/courses/18-433-combinatorial-optimization-fall-2003/resources/a1/}{官方作业1}，PDF第1页；编者独立解答，非官方答案}
\begin{solution}
把边按权非递增排序后，按题2的一次扫描规则选择，时间 $O(n+m\log(m+1))$。设贪心匹配为 $M$，最优为 $M^*$。每条最优边 $e$ 若也被贪心选中，就记账给自己；否则它被扫描时已有端点被此前选入的某条边 $f$ 占用，记账给这样的 $f$。排序保证 $w(e)\le w(f)$。

每条 $f\in M$ 至多接受两条最优边的记账，因为最优匹配在 $f$ 的每个端点至多有一条边。若 $f$ 本身属于 $M^*$，则不会再有其他最优边与它相交。于是
\[
 w(M^*)=\sum_{e\in M^*}w(e)\le\sum_{f\in M}2w(f)=2w(M).
\]
非负性使未收到记账的边也可以纳入右边而不降低总和。其必要性可见于仅有一条权为 $-1$ 的边：强制选到极大的算法输出 $-1$，空匹配的最优值为0。另一种有效修订是遇到非正权即停止，但那属于更改原选择规则。
\end{solution}

\begin{exercise}\label{mit:18433-a1-q4}
\textbf{18.433 Fall 2003，作业1，题4.}
称一个图性质具有“良好刻画”，若它成立和不成立时都能给出长度关于顶点数 $n$ 为多项式、并可有效核查的证明。课堂上二分图有完美匹配的性质具有这种刻画。判断并证明以下性质是否也有良好刻画：
\begin{enumerate}[label=(\alph*)]
\item 图 $G$ 连通。
\item $G$ 是 Euler 图，即存在恰经过每条边一次的闭合巡回。
\item $G$ 可用两种颜色给顶点染色，使每条边两端异色。
\end{enumerate}
编者约定：$n\ge1$；Euler 巡回不必访问孤立顶点，无边图允许长度0的闭行走。
\end{exercise}
\sourceof{\href{https://ocw.mit.edu/courses/18-433-combinatorial-optimization-fall-2003/resources/a1/}{官方作业1}，PDF第1页；编者独立解答，非官方答案}
\begin{solution}
三种性质都有良好刻画。这里不仅给肯定情形的证书，也给否定情形的证书。

(a) 肯定证书是一棵覆盖全部顶点的生成树：检查它的 $n-1$ 条边属于 $G$、连通且无圈即可。否定证书是非空真子集 $S\subset V$，检查没有边连接 $S$ 与 $V\setminus S$。不连通时取一个连通分量便得到这种 $S$；连通时逐步加入连接新顶点的边便得到生成树。

(b) 肯定证书直接列出 Euler 闭合巡回。它有 $m$ 条边，核查相邻边首尾相接、末点为起点、每条原边恰出现一次即可。否定证书或者是一个奇度顶点，或者是两个含边的连通分量之间的分割。前者不可能在闭行走中把每次进入与离开配对，后者使一条行走无法经过两分量的边。还须证明这些障碍穷尽：若所有非零度顶点处于一个连通分量且度均为偶数，从其中一点不断走未用边，除起点外不可能无边可走，因每次到达所消耗的度为奇数，尚有未用边离开；所以得到闭合游走。若仍有未用边，连通性保证已有游走某点关联未用边，从该点再构造闭合游走并插入。每次消耗至少一条新边，最后得到所需巡回。

(c) 肯定证书是两侧顶点名单，逐边查两端异侧。否定证书是一个奇圈，查其各边属于图及长度为奇即可。题1已经证明，不能二染色恰好会出现奇圈。

简单图有 $m\le\binom n2$，故上述顶点名单、树、巡回和圈的长度均关于 $n$ 为多项式；每项核查可用图搜索或扫描边表完成。
\end{solution}

\begin{exercise}\label{mit:18433-a1-q5}
\textbf{18.433 Fall 2003，作业1，题5.}
二分图的 Hungarian 匹配算法用交错森林寻找增广路。其分阶段改进版在每一阶段同时沿一组极大的彼此不相交增广路增广。证明：每一阶段后，最短增广路长度严格增加。

编者补条件：本阶段选取的均为当前匹配的\textbf{最短}增广路，“不相交”指顶点不相交，“极大”指无法再加入一条与已选路顶点不相交的同长度增广路。原题本页未明确最短限定；任意长度路组的极大性不足以支持结论。阶段后若已无增广路，最短长度记为 $+\infty$。
\end{exercise}
\sourceof{\href{https://ocw.mit.edu/courses/18-433-combinatorial-optimization-fall-2003/resources/a1/}{官方作业1}，PDF第1--2页；编者独立解答，非官方答案}
\begin{solution}
设二分为 $A\sqcup B$，阶段前匹配为 $M$。建立有向辅助图：非匹配边定向为 $A\to B$，匹配边为 $B\to A$；加虚源 $s$ 到每个自由 $A$ 点、每个自由 $B$ 点到虚汇 $t$ 的弧。$s$--$t$ 路恰对应增广路，长度比原路多2。记阶段前从 $s$ 的距离为 $d$，$d(t)=L$。所选各最短路的每条弧 $u\to v$ 都满足 $d(v)=d(u)+1$，否则换用更短前缀会得到更短源汇行走，去圈后仍更短，矛盾。

同时翻转顶点不相交的增广路，得到匹配 $M'$。新辅助图删除了被匹配的自由端点的源、汇弧；其余改变是所选路上的图边反向。每条旧弧满足 $d(v)\le d(u)+1$；新反向弧 $v\to u$ 则满足 $d(u)=d(v)-1$，也满足距离三角不等式。因此沿新图任何从 $s$ 到 $t$ 的路累加不等式，长度至少为 $L$。旧时不可达的顶点也不会因新增反向弧首次可达，因为这些反向弧两端都在原最短路上。

若新图仍有长度 $L$ 的源汇路 $Q$，其每一步都必须令距离增加1，不能用任何新反向弧，因为那种弧令旧距离减少1。再说明 $Q$ 不能与任何所选增广路共用图顶点。所选路上的每点在 $M'$ 中都已匹配；若它是 $Q$ 的端点，就违反端点自由性；若它是内部点，$Q$ 必须使用它唯一的 $M'$ 匹配边。这条匹配边是所选路原先非匹配而现已翻转的边，沿 $Q$ 的方向正是新反向弧，也不可能。

所以 $Q$ 对应一条阶段前的同长度增广路，并与全部已选路顶点不相交，可再加入所选组，违反极大性。故新最短长度严格大于旧长度。

最短限定确为必要。取两侧 $A=\{a_0,a_1,a_2,a_3\}$、$B=\{b_0,b_1,b_2,b_3\}$，原匹配 $M=\{a_1b_1,a_2b_2\}$，再加边 $a_0b_1,a_1b_2,a_2b_0,a_3b_2,a_1b_3$。原最短增广路长3，例如 $a_0,b_1,a_1,b_3$。若只沿长5路 $a_0,b_1,a_1,b_2,a_2,b_0$ 增广，这一单路组在全部增广路中已是顶点不相交意义的极大组：未使用顶点只有 $a_3,b_3$，两者无边。新匹配却仍有长3增广路 $a_3,b_2,a_1,b_3$，最短长度没有增加。
\end{solution}

\begin{exercise}\label{mit:18433-a1-q6}
\textbf{18.433 Fall 2003，作业1，题6.}
设 $G=(V,E)$ 无孤立顶点。顶点覆盖与每条边至少有一端相交；独立集内任意两点不相邻；边覆盖是一组边，使每个顶点至少关联其中一条边。记 $\alpha(G)$ 为最大独立集大小，$\rho(G)$ 为最小边覆盖大小，$\tau(G)$ 为最小顶点覆盖大小，$\nu(G)$ 为最大匹配大小。证明
\[
 \alpha(G)+\tau(G)=|V|=\nu(G)+\rho(G).
\]
原提示：第一式找独立集与覆盖的对应；第二式分别从最大匹配构造小边覆盖、从最小边覆盖构造大匹配。
\end{exercise}
\sourceof{\href{https://ocw.mit.edu/courses/18-433-combinatorial-optimization-fall-2003/resources/a1/}{官方作业1}，PDF第2页；编者独立解答，非官方答案}
\begin{solution}
集合 $I$ 独立当且仅当补集 $V\setminus I$ 是顶点覆盖。因此取最大、最小分别得到 $\tau=|V|-\alpha$。

取大小 $\nu$ 的最大匹配。其未覆盖顶点共有 $|V|-2\nu$ 个，彼此不相邻，否则还能加入一条边。每个未覆盖点因无孤立点而有邻边，选一条加到匹配中，得到边覆盖，最多使用
$\nu+(|V|-2\nu)=|V|-\nu$ 条边。所以 $\rho\le|V|-\nu$。

反向取最小边覆盖 $F$。如果 $F$ 中某条边两端在 $(V,F)$ 中的度都至少2，删去它后仍覆盖所有点，与最小性矛盾。因此 $F$ 的每条边至少有一个度1的端点。每个分量必为星：若含圈，圈边两端度至少2；若一个分量有两个度至少2的顶点，连接它们的最短路上也会有两端度至少2的边，均矛盾。每个星有 $v_i$ 点、$v_i-1$ 条边；每个分量选一条边，彼此组成匹配，其大小为分量数
\[
 c=\sum_i\bigl(v_i-(v_i-1)\bigr)=|V|-|F|=|V|-\rho.
\]
所以 $\nu\ge|V|-\rho$，与前一个不等式合起来即得第二式。空图的四个参数都为0，也符合结论。
\end{solution}

\begin{exercise}\label{mit:18433-a1-q7}
\textbf{18.433 Fall 2003，作业1，题7.}
图中独立集是内部没有边的顶点集合。
\begin{enumerate}[label=(\alph*)]
\item 给出一个整数规划，其可行解恰对应 $G=(V,E)$ 的独立集。
\item 放松整数限制得到一个多面体。证明若 $G$ 为二分图，这个多面体的全部顶点都是整数点。
\end{enumerate}
\end{exercise}
\sourceof{\href{https://ocw.mit.edu/courses/18-433-combinatorial-optimization-fall-2003/resources/a1/}{官方作业1}，PDF第2页；编者独立解答，非官方答案}
\begin{solution}
(a) 为每点设变量 $x_v$，$x_v=1$ 表示选它。约束为
\[
 x_u+x_v\le1\quad(uv\in E),\qquad x_v\in\{0,1\}\quad(v\in V).
\]
若要求最大独立集，再加目标 $\max\sum_v x_v$；仅描述所有独立集时不需目标。边约束正好排除同时选择相邻两点。

(b) 松弛为
\[
 P=\{x\in\R^V:0\le x_v\le1,\ x_u+x_v\le1\ (uv\in E)\}.
\]
这里多面体指有限个线性半空间的交；它的顶点也叫极点，意思是不能写成两个不同可行点的非平凡凸组合。本题直接证明极点整性，不依赖线性规划对偶或整矩阵定理。

假设 $x\in P$ 有分数坐标。令 $F=\{v:0<x_v<1\}$，在 $F$ 上只保留紧边，即满足 $x_u+x_v=1$ 的边，取其中任一连通分量 $C$，单个孤点也算分量。原图二分为 $A\sqcup B$，定义方向 $h$ 在 $C\cap A$ 上为1，在 $C\cap B$ 上为 $-1$，其余点为0。$h$ 非零。

与 $C$ 有关的每条紧边两端都在 $C$ 中，故其两端方向和为0。说明不会遗漏一端在 $C$、另一端为整数坐标的紧边：若另一端为0，紧性迫使本端为1；若另一端为1，本端为0，均与本端分数矛盾。一端在 $C$ 而另一分数点不在 $C$ 的紧边也不可能，因为它会把两点连入同一分量。

因此所有可能改变的边约束都有正余量；$C$ 中的上下界也都有正余量。约束有限，可取足够小的 $\varepsilon>0$，使 $x+\varepsilon h$、$x-\varepsilon h$ 都满足上下界和这些非紧约束，而紧约束保持不变。这两个可行点不同，且 $x$ 是它们的中点，所以 $x$ 不是极点。矛盾证明全部极点均为0--1点。上界 $x_v\le1$ 对孤立点也必须保留，不能因某点没有边约束而漏写。
\end{solution}

\par\addvspace{7pt}\Needspace{4\baselineskip}\noindent\textbf{作业2：流理论}

\begin{exercise}\label{mit:18433-a2-q1}
\textbf{18.433 Fall 2003，作业2，题1.}
设有向图 $G$ 有不同端点 $s,t$。两条从 $s$ 到 $t$ 的有向路，若除端点外没有共同顶点，称为内部顶点不相交。证明最多内部顶点不相交路的条数，等于删去顶点后使 $s$ 到 $t$ 不再有路的最少顶点数。

编者补条件：$G$ 无弧 $s\to t$，只能删除 $V\setminus\{s,t\}$ 中的顶点。原页未明说这两点；若已有直接弧，删任何内部点都不能切断它，原等式不能按此含义成立。
\end{exercise}
\sourceof{\href{https://ocw.mit.edu/courses/18-433-combinatorial-optimization-fall-2003/resources/a2/}{官方作业2}，PDF第1页；编者独立解答，非官方答案}
\begin{solution}
记 $n=|V|$，取整数 $B=n+1$。把每个内部点 $v$ 拆成 $v^-,v^+$，加容量1的弧 $v^-\to v^+$；每条原弧 $u\to v$ 变成从 $u$ 的出端到 $v$ 的入端的容量 $B$ 弧，$s,t$ 不拆。进入源和离开汇的原弧可以删除，它们不可能用于简单源汇路。无直接弧保证删除全部 $n-2$ 个内部点就能分离，因此拆点网络有一个容量至多 $n-2$ 的割。

任取原图内部顶点分离集 $X$，删去对应拆点弧后汇不可达；取剩余网络从源可达集，其出割只包含被删的那些单位弧，容量至多 $|X|$。反向，拆点网络最小割容量小于 $B$，不含任何原图转来的大容量弧，只包含若干单位拆点弧。它们对应顶点集 $X$；若原图还有避开 $X$ 的源汇路，转成拆点路就会跨该割而不使用任何被割单位弧，矛盾。因此 $X$ 是分离集，且大小等于该割容量。

整数最大流最小割定理给一个整数最大流。先从正流中消去有向圈，再反复抽取一单位源汇路流；单位拆点弧容量1使不同单位路不能经过同一内部点。于是最大流值给出这么多内部顶点不相交原图路。反过来，$k$ 条这种路转成单位流相加仍满足所有拆点容量，大容量 $B$ 也足以承接，故流值至少 $k$。两向对应和前段的割对应一起证明所求等式。无路时两者皆为0。
\end{solution}

\begin{exercise}\label{mit:18433-a2-q2}
\textbf{18.433 Fall 2003，作业2，题2.}
有向图的边有容量，指定不同顶点 $s,t$。一条有向源汇路的容量是其边容量的最小值。给出有效算法，寻找容量最大的源汇路。

容量按网络的通常含义为非负数；若无源汇路，算法报告不存在，而不是输出一条容量为0的虚构路径。
\end{exercise}
\sourceof{\href{https://ocw.mit.edu/courses/18-433-combinatorial-optimization-fall-2003/resources/a2/}{官方作业2}，PDF第1页；编者独立解答，非官方答案}
\begin{solution}
这是最大瓶颈路，不是边容量之和最大的路。对每点维护当前最好瓶颈 $b(v)$，初始 $b(s)=+\infty$，其他点为 $-\infty$，后者表示尚无已知路；记录前驱。每次取尚未确定而 $b$ 最大的点 $u$，确定它，逐条松弛出弧 $u\to v$：若
\[
 \min\{b(u),c(u,v)\}>b(v),
\]
就用该值更新 $b(v)$ 和前驱。到确定 $t$ 时沿前驱恢复路；若剩余键全为 $-\infty$ 而汇尚未确定，则无路。这个初值还正确区分了真实零容量路和不存在路径。

证明确定点的键恰为最优瓶颈。已有键总由真实路径产生，所以不超过最优。假设某个即将确定的 $u$ 还有瓶颈更大的源到 $u$ 路，取该路第一个未确定点 $v$，其前驱 $w$ 已确定。归纳假设保证 $b(w)$ 至少是该路前缀瓶颈；确定 $w$ 时松弛 $wv$，使 $b(v)$ 至少是整条路的瓶颈，因而严格大于 $b(u)$。这与当前取最大键矛盾。初始源显然正确，归纳完成。

用邻接表和带位置索引的二叉最大堆，确定点至多 $n$ 次、更新至多 $m$ 次，总时间 $O((n+m)\log n)$，额外存储 $O(n+m)$。这里是每点一个堆项的实现，平行弧不会把对数因子偷偷改成堆中边项数的对数。
\end{solution}

\begin{exercise}\label{mit:18433-a2-q3}
\textbf{18.433 Fall 2003，作业2，题3.}
称一个割在最小割的 $k$ 倍以内，若其边数至多为最小割边数的 $k$ 倍。设 $2k$ 为整数。证明：任一无向图中，这种割的数目严格少于 $n^{2k}$。

编者补条件与计数约定：$n\ge2$，图连通，$k\ge1$，割指非平凡顶点二分且两侧交换不重复计数。允许平行边，按其重数计边，删去自环。原题未写连通性；最小割为0的非连通图不能直接使用下面的结论及收缩估计。
\end{exercise}
\sourceof{\href{https://ocw.mit.edu/courses/18-433-combinatorial-optimization-fall-2003/resources/a2/}{官方作业2}，PDF第1页；编者独立解答，非官方答案}
\begin{solution}
设最小割大小为 $\lambda\ge1$，令 $r=2k\ge2$ 为整数。先处理 $n\ge r$。均匀随机缩边，直到剩 $r$ 个超级顶点，然后在它们的 $2^{r-1}-1$ 个无序非平凡二分中均匀选一个，输出对应原图割。

固定一个目标割 $C$，大小 $|C|\le k\lambda=r\lambda/2$。条件于此前没有缩掉 $C$ 的边，它仍为当前图的割。当前任何割都是原图割，所以每个超级顶点度至少 $\lambda$；剩 $j$ 点时边数至少 $j\lambda/2$。下一步碰到目标割的概率至多 $r/j$。故生存到 $r$ 点的概率至少
\[
 \prod_{j=n}^{r+1}\left(1-\frac rj\right)
 =\frac{r!(n-r)!}{n!}=\binom nr^{-1}.
\]
生存后目标二分仍非平凡，最终均匀二分恰好输出它的条件概率为 $1/(2^{r-1}-1)$。于是每个目标割的输出概率至少为
$1/[\binom nr(2^{r-1}-1)]$。不同输出互斥，若目标割共有 $N$ 个，则
\[
 N\le\binom nr(2^{r-1}-1)
 \le\frac{n^r}{r!}(2^{r-1}-1)<n^r.
\]
最后严格不等式来自 $2^{r-1}-1<r!$：$r=2$ 时为 $1<2$；此后 $r!$ 每次至少乘3，而 $2^{r-1}-1$ 增长至多为前项的3倍。

若 $n<r$，全部非平凡无序割也只有 $2^{n-1}-1$ 个，而 $2^{n-1}\le n^{n-1}<n^r$。两种情况均得 $N<n^{2k}$。这一论证通过一个输出分布同时约束全部目标割，没有把对不同割的生存事件误当成彼此独立。

连通条件也是实质补充：8个孤立点组成的图，最小割为0，全部 $2^7-1=127$ 个无序非平凡割都为0；取 $k=1$ 时，原题未补条件的断言会要求 $127<8^2=64$，并不成立。
\end{solution}

\begin{exercise}\label{mit:18433-a2-q4}
\textbf{18.433 Fall 2003，作业2，题4.}
设 $P$ 遍历图中两不同顶点 $s,t$ 间的路径，$C$ 遍历分离 $s,t$ 的割，$c_e$ 是边容量。证明
\[
 \max_P\min_{e\in P}c_e=\min_C\max_{e\in C}c_e.
\]
编者约定：容量非负；无路径时左边为0，空割的最大容量为0。有向版本的割只计从含 $s$ 的一侧指向含 $t$ 的一侧的弧，无向版本则计全部跨边。
\end{exercise}
\sourceof{\href{https://ocw.mit.edu/courses/18-433-combinatorial-optimization-fall-2003/resources/a2/}{官方作业2}，PDF第1页；编者独立解答，非官方答案}
\begin{solution}
记左边为 $b$。每条源汇路都必须跨每个源汇割，所以对任何 $P,C$，任选其一条跨割边 $e$ 就有
\[
 \min_{f\in P}c_f\le c_e\le\max_{f\in C}c_f.
\]
对路径取最大、对割取最小，得左边不超过右边。

为取反向界，仅保留容量严格大于 $b$ 的边，令 $S$ 为从 $s$ 可达点集。$t\notin S$：否则该子图的一条源汇路所有边容量都大于 $b$，违背 $b$ 的定义。$S$ 的每条出割边容量都至多 $b$，否则其终点也会可达。故这个割满足最大边容量至多 $b$，右边不超过左边。有无路径的边界按题面约定同样成立。

右边是割中\textbf{最大一条边的容量}，并非跨割容量之和；因此不能把本式直接当成最大流最小割定理。题2的最大瓶颈路算法能计算左边，阈值可达集又给出了相等的割证书。
\end{solution}

\begin{exercise}\label{mit:18433-a2-q5}
\textbf{18.433 Fall 2003，作业2，题5.}
设最大流增广算法每次都选择所含反向弧最少的增广路。这里反向弧指要在其原弧上撤回流量的残量弧。给出总增广次数的上界。

编者明确模型：从零流开始，每次沿所选路增广其正瓶颈量；前向、后向残量弧保留原弧编号。容量可以是非负实数，数学证明使用精确运算，不增添整数容量假设。
\end{exercise}
\sourceof{\href{https://ocw.mit.edu/courses/18-433-combinatorial-optimization-fall-2003/resources/a2/}{官方作业2}，PDF第1页；编者独立解答，非官方答案}
\begin{solution}
给每个前向残量身份成本0、后向身份成本1。题中的选择正是这个0--1成本下的最短源汇路；它与最少总弧数的 Edmonds--Karp 选择不同。设 $d(v)$ 为当前从源到 $v$ 的最少成本，不可达为 $+\infty$。非负成本允许去掉圈而不增加成本，所以有限距离是 $0,1,\ldots,n-1$ 中的整数，所选路的每条弧 $u\to v$ 均为紧弧：
\[
 d(v)=d(u)+w(u,v).
\]
若某前缀不是最短，换成较短前缀接后缀，去圈后会给更小成本的源汇路，矛盾。

先证距离不减。旧残量弧都满足 $d(v)\le d(u)+w(u,v)$。增广后新出现的残量身份只可能是所用弧的相反身份；两种身份成本之和为1。若旧路用 $u\to v$，成本为 $w\in\{0,1\}$，新反向 $v\to u$ 的成本为 $1-w$，则
\[
 d(u)=d(v)-w\le d(v)+(1-w).
\]
所以旧距离对所有新弧仍满足三角不等式，沿任何新源到点的路累加，就得新距离至少是旧距离。旧时不可达的顶点也不能由新增身份首次到达，因为新增身份的两端都已在旧增广路上可达。

每次增广至少使一条残量身份达到瓶颈而饱和，称它为本次的临界身份。某身份 $u\to v$ 成本为 $w$，第一次临界时记尾点距离 $a=d(u)$，则 $d(v)=a+w$。它消失后若想再次临界，先须在某次中间增广中使用相反身份 $v\to u$，使其残量重新为正。那次最短路的距离 $d'$ 满足
\[
 d'(u)=d'(v)+(1-w)\ge d(v)+(1-w)=a+1.
\]
距离此后仍不减。因此同一身份每次重新临界，尾点整数距离至少增加1；尾点在当次源汇路上，有限距离至多 $n-1$，故同一身份至多临界 $n$ 次。原弧有 $m$ 条，残量身份至多 $2m$ 个，把每次增广记给其中一个临界身份，得总次数至多 $2nm$。

用双端队列的0--1最短路算法，一次搜索需 $O(n+m)$ 时间：成本0的改进放队首，成本1的改进放队尾，按非递减距离处理并扫描出弧。这与二层桶的 Dijkstra 等价，每个顶点最终确定一次。连同路径更新，总时间为 $O(nm(n+m))$ 次基本运算；终止后无残量路，可达集给最大流最小割证书。证明对精确实容量成立；有理输入采用二进制表示时，所有流量可用原容量分母的公共倍数表示，数值位长亦为输入长度的多项式。浮点舍入的停机判定不属于这个精确算法结论。
\end{solution}

\par\addvspace{7pt}\Needspace{4\baselineskip}\noindent\textbf{作业3：凸规划中的图优化题}

\begin{exercise}\label{mit:18433-a3-q6}
\textbf{18.433 Fall 2003，作业3，题6.}
完全图 $G=(V,E)$ 的每条边 $ij$ 有正长度 $w_{ij}$；旅行商问题要求寻找总长度最小的 Hamilton 圈。
\begin{enumerate}[label=(\alph*)]
\item 给出解决旅行商问题的整数线性规划。
\item 放松整数限制得到线性规划；通过设计分离预言机，证明该松弛可以在多项式时间求解。
\end{enumerate}
编者明确：$n=|V|\ge3$。第(b)问的位复杂度以正有理长度的二进制编码为输入模型；原题未规定数的表示方式。
\end{exercise}
\sourceof{\href{https://ocw.mit.edu/courses/18-433-combinatorial-optimization-fall-2003/resources/a3/}{官方作业3}，PDF第2页；编者独立解答，非官方答案}
\begin{solution}
(a) 用 $x_e\in\{0,1\}$ 表示是否选边，$\delta(S)$ 表示恰有一个端点在 $S$ 的边集，整数规划为
\begin{align*}
 \min\quad&\sum_{e\in E}w_e x_e,\\
 \text{满足}\quad&\sum_{e\in\delta(\{v\})}x_e=2&& (v\in V),\\
 &\sum_{e\in\delta(S)}x_e\ge2&& (\varnothing\ne S\subsetneq V),\\
 &x_e\in\{0,1\}&& (e\in E).
\end{align*}
Hamilton 圈显然每点度2，且每个非平凡割至少跨出两次，因此给可行整数点。反过来，0--1点选择的子图每点度2，每个分量为圈；若不连通，取一个分量顶点集 $S$，其出割和为0，违反割约束。所以子图连通，是一个 Hamilton 圈。只有度等式会允许若干互不相交的子圈，必须保留割约束。

(b) 改为 $0\le x_e\le1$，保留度等式和割不等式，得到有界有理多面体 $P$。以下先交代此问调用的优化工具，再实际给出分离算法。

\textbf{基础优化算法定理（本节明确调用）.}
对一个由有理线性等式、不等式描述的多面体，若维数、各行系数的编码长度和已知外界的编码长度有多项式界，且存在强分离预言机：给有理点时，要么正确确认其可行，要么输出一条被它违反、而对全部可行点有效的有理线性约束；预言机的运行时间及输出位长对输入编码为多项式，那么可用有理椭球法在多项式时间进行精确线性优化，返回最优值和一个最优有理点，或报告不可行。给出的等式可先用有理 Gaussian 消元转到其仿射解空间；多面体低维不应被误判为空。基础定理包含有限精度处理和由有理顶点编码界恢复精确解的部分。本节把它作为原课线性规划背景调用，不声称在这里证明整套有理椭球算法。课堂工具可参见18.433的\href{https://ocw.mit.edu/courses/18-433-combinatorial-optimization-fall-2003/resources/l1617/}{椭球法讲义}和\href{https://ocw.mit.edu/courses/18-433-combinatorial-optimization-fall-2003/resources/l18/}{分离预言机讲义}。这里准确陈述所调用的有理算法版本，不把短课堂稿当成其有限精度部分的完整证明；下面的图论分离算法及条件核对则完整证明。

先检查输入 $x$ 的每条界约束 $0\le x_e\le1$，再检查每点的度和是否为2。发现错误时，直接输出相应约束；若等式的左边大于或小于2，分别用同方向的 $\le2$ 或 $\ge2$ 给分离半空间。这样的半空间包含 $P$ 并严格排除所给点。此步骤为 $O(n+m)$ 次有理运算，$m=\binom n2$。

通过这些检查后，以非负数 $x_e$ 为边容量求一个全局最小割。可固定根 $r$，对每个 $t\ne r$ 求 $r$--$t$ 最小割并取最小者：任一无序非平凡割都有不含 $r$ 的一侧，其中取 $t$，它就是一个 $r$--$t$ 割；故上述最小值不大于任意割容量，同时每次返回的割也是全局候选割，所以恰为全局最小值。无向边变为两个相反方向的容量 $x_e$ 原弧时，某一出割恰计每条无向跨边一次，容量保持不变。可用第12章已证明的 Edmonds--Karp 算法求这些有理最大流和割。

若最小值小于2，输出它的顶点集 $S$ 及约束 $\sum_{e\in\delta(S)}x_e\ge2$，这条约束对整个 $P$ 有效且严格排除当前点。若最小值至少2，则全部指数多个割约束均成立，可以确认当前点属于 $P$。预言机无须把全部割先枚举出来。

$n-1$ 次最大流，每次的运算界为 $O(nm(n+m))$，故此预言机可用 $O(n^2m(n+m))$ 次基本运算，所得分离行只有0、1系数和常数2，编码长度为多项式。输入的容量为有理数，精确流运算的位长如本节作业2第5题所述为多项式；没有把实数单位成本模型误充二进制时间界。

最后核对基础优化定理的低维接口。$P\subseteq[0,1]^m$ 给已知外界，Hamilton 圈给非空性。度等式矩阵在完全图 $n\ge3$ 上秩为 $n$：若其行的线性组合为0，系数 $a_v$ 必满足 $a_u+a_v=0$ 对每一对不同点成立；取三个点即推出它们系数全为0，其他点亦为0。$n=3$ 时只有一个可行点，三条边全为1，直接输出即可。

$n\ge4$ 时取 $z_e=2/(n-1)$。它满足度等式，且在所有边界的内部；对 $2\le |S|\le n-2$，
\[
 \sum_{e\in\delta(S)}z_e
 =\frac{2|S|(n-|S|)}{n-1}
 =2+\frac{2(|S|-1)(n-|S|-1)}{n-1}>2.
\]
单点割和其补集的割就是度等式。因此除这些必等的行以外均有正余量，$P$ 的仿射包正是度等式的解空间。用有理消元选自由坐标后，维数为 $m-n$，代换系数位长由0--1矩阵的子式界控制，为多项式；预言机输出的行按同一代换投到该空间，仍合法且位长有多项式界。算法得到自由坐标的最优有理点后，按消元关系恢复全部 $x_e$；目标也按同一关系代换，并保留常数项。这说明求解及恢复接口完整满足所调用定理。

所得通常是分数边向量，解决的是线性规划松弛；若它恰为0--1，才可以按第(a)问恢复 Hamilton 圈，不能把一般分数最优点当成原旅行商问题的精确解。
\end{solution}

\par\addvspace{7pt}\Needspace{4\baselineskip}\noindent\textbf{作业4：近似算法中的图论题}

\begin{exercise}\label{mit:18433-a4-q2}
\textbf{18.433 Fall 2003，作业4，题2.}
无向图 $G=(V,E)$ 的生成子图 $G'=(V,E')$ 若任意两点间至少有两条边不相交路径，称为二边连通生成子图。
\begin{enumerate}[label=(\alph*)]
\item 证明寻找边数最少的二边连通生成子图是 NP-hard 的。
\item 给出有效的二近似算法，即输出边数至多为最优值两倍的可行生成子图。
\end{enumerate}
编者明确：采用有限简单图、$n\ge3$；若原图不二边连通，先报告无可行解。
\end{exercise}
\sourceof{\href{https://ocw.mit.edu/courses/18-433-combinatorial-optimization-fall-2003/resources/a4/}{官方作业4}，PDF第1页；编者独立解答，非官方答案}
\begin{solution}
(a) 使用无向 Hamilton 圈问题的 NP-complete 性作为复杂度背景。归约把 $n$ 点图原样交给判定问题：“是否有至多 $n$ 条边的二边连通生成子图？”若有 Hamilton 圈，它的 $n$ 条边就是肯定证书。反过来，二边连通子图每点度至少2，握手引理给边数至少 $n$。若又至多 $n$ 条，就每点度恰为2，且子图连通，只能是一个 Hamilton 圈。因此两判定答案相同，归约只复制输入，是多项式时间，最小化问题由此 NP-hard。即使限定原图已二边连通也不改变困难性：先检测任意 Hamilton 输入是否连通且无桥；若不满足，它必无 Hamilton 圈，可把它映到两三角形共一个顶点的固定五点图。该图二边连通，却因另外四点的度2强制保留全部六边，不存在至多五边的可行子图；其余输入原样保留。检查生成性、连通性与是否有桥也是多项式时间，故相应判定问题还属于 NP。

(b) 先解释可行性检查。连通无向图二边连通当且仅当没有桥：桥阻断任意两点的第二条边不相交路；没有桥意味着每个非平凡边割至少两条边，由边版 Menger 得任意两点至少两条边不相交路。可用一次 DFS 检测桥，或逐边删除并搜索也仍为多项式时间。

对可行输入，下面实际构造一个边数不超过 $2n-3$ 的二边连通生成子图。先找任一圈 $H$，设它有 $\ell\ge3$ 个顶点。保持 $H$ 连通且无桥。若尚未包含全部顶点，取 $G-V(H)$ 的一个连通分量 $C$。它到 $H$ 至少有两条不同的连接边，否则唯一连接边是原图的桥。选两条连接边 $ux,vy$，其中 $u,v\in V(H)$、$x,y\in C$，允许 $u=v$ 或 $x=y$；在 $C$ 中找一条简单 $x$--$y$ 路，$x=y$ 时取零长路。两连接边加这条路形成一条内部顶点都新、两端在 $H$ 的路；两端相同时是闭合的耳。

把这条耳加入 $H$。旧边仍位于旧圈上。新边若两端不同，就与旧 $H$ 中一条端点间路径组成圈；两端相同则耳自身成圈。因此每条新边也位于圈上，子图仍连通无桥。每轮至少增加一个新顶点，所以最终得到生成子图。

若一轮加 $p$ 个新顶点，正好加 $p+1$ 条边。设加耳 $q$ 轮，则
\[
 |E(H)|=\ell+\sum_{i=1}^q(p_i+1)=n+q,
 \qquad q\le\sum_i p_i=n-\ell.
\]
故边数至多 $2n-\ell\le2n-3$。任何可行最优子图每点度至少2，边数 $\mathrm{OPT}\ge n$，从而 $|E(H)|<2\,\mathrm{OPT}$，当然满足二近似要求。每轮通过搜索找分量、连接边和内部路需 $O(n+m)$ 时间，至多 $n$ 轮，总时间 $O(n(n+m))$。这一证明允许闭合耳，不错误地把二边连通图当成二顶点连通图。
\end{solution}

\begin{exercise}\label{mit:18433-a4-q3}
\textbf{18.433 Fall 2003，作业4，题3.}
有向图中的无圈子图是一组不包含任何有向圈的边 $E'\subseteq E$。
\begin{enumerate}[label=(\alph*)]
\item 给一个多项式时间算法，使输出无圈子图的边数至少为最优值的一半。
\item 以每条边一个变量，给出最大基数无圈子图的整数规划。
\item 放松整数限制得到 LP；它的整性间隙可以多大？本题规定比值为 LP 最优值除以整数最优值。
\item 给出多项式时间分离预言机，从而求解该 LP。
\end{enumerate}
编者约定：自环不能入无圈子图，先删除；平行弧和两个同端点反向弧各自保留身份。有边时讨论整性间隙；删自环后无边的零最优情形不取 $0/0$。
\end{exercise}
\sourceof{\href{https://ocw.mit.edu/courses/18-433-combinatorial-optimization-fall-2003/resources/a4/}{官方作业4}，PDF第1页；编者独立解答，非官方答案}
\begin{solution}
(a) 任取顶点全序，把指向序号较大点的弧记为 $F$，反向弧记为 $B$。沿 $F$ 的路序号严格增加，沿 $B$ 的路严格减少，所以两者都无圈。输出较大者，大小至少 $m/2\ge\mathrm{OPT}/2$。记录序号后扫描每弧一次，时间 $O(n+m)$。又因为任何无圈子图存在拓扑序，它的全部弧都在该序的前向弧中，故
\[
 \mathrm{OPT}=\max_{\pi}\bigl|F_\pi\bigr|.
\]
拓扑序的存在可直接由反复删除入度0点证明；若剩余子图没有这种点，逆着入弧追踪会遇到重复点并形成有向圈。

(b) 令 $x_e$ 表示选弧，最大化 $\sum_e x_e$，约束为
\[
 \sum_{e\in C}x_e\le |C|-1\quad(\text{每个有向简单圈 }C),
 \qquad x_e\in\{0,1\}.
\]
若选弧包含有向圈，取其中一个简单圈，该行左边为 $|C|$，不满足；若无有向圈，每个圈至少漏一条弧，所有行均满足。因此整数点恰对应无圈子图。两个同端点反向弧形成长度2的圈，也要列入约束，不因顶点少而漏掉。

(c) 松弛为 $0\le x_e\le1$。它的整性间隙上确界为2：一方面 LP 值至多 $m$，第(a)问给 $\mathrm{OPT}\ge m/2$，所以比值不超过2。下面构造趋近2的实例，补足只有上界尚不能说明间隙规模的缺口。

固定整数 $g\ge3$ 和很小的 $0<\delta<1/10$。在 $n$ 点的每对无序点之间，独立以概率 $p=n^{-1+1/(2g)}$ 放一条边；若放边，独立等概率选择两个方向之一。记 $N=\binom n2$，随机弧数为 $M$，其均值为 $pN$、方差至多 $pN$。由方差的 Markov 不等式，
\[
 \Pr\bigl(|M-pN|>\delta pN\bigr)\le\frac1{\delta^2pN}\longrightarrow0.
\]
对固定次序 $\pi$，令每对点贡献 $Y=1$（前向）、$-1$（后向）或0（无边），则 $F_\pi-M/2=\tfrac12\sum Y$。$Y$ 的指数矩为
\[
 \E e^{tY}=1-p+p\cosh t
 \le\exp\bigl(p(\cosh t-1)\bigr)\le e^{pt^2}
 \quad(0\le t\le1).
\]
最后一步可从幂级数核查：$\cosh t-1=\sum_{j\ge1}t^{2j}/(2j)!\le t^2(\cosh1-1)<t^2$。独立性与 Markov 不等式，取 $t=\delta/2$，给
\[
 \Pr\!\left(\sum Y>\delta pN\right)
 \le\exp(-t\delta pN+pt^2N)
 =\exp(-\delta^2pN/4).
\]
次序至多 $n!\le n^n$ 个，故至少有一次序超过界的概率至多
$\exp(n\log n-\delta^2pN/4)\to0$，因为 $pN$ 的阶为 $n^{1+1/(2g)}$。所以以趋于1的概率，\textbf{全部}次序同时满足 $F_\pi\le M/2+\delta pN/2$。

再令 $X$ 为长度小于 $g$ 的有向简单圈数。长度2的圈不会出现；对 $3\le j<g$，候选圈数至多 $n^j$，每条所需定向弧同时出现的概率为 $(p/2)^j$，因此
\[
 \E X\le\sum_{j=3}^{g-1}(np/2)^j
 \le g\,n^{(g-1)/(2g)}.
\]
当 $g=3$ 时该和为空。由 Markov，$\Pr(X>n)\le\E X/n\to0$。综合三项概率界，当 $n$ 足够大时存在一张图，同时满足弧数接近期望、所有次序的前向弧数受界及 $X\le n$。

从这种图的每个短圈任选一条弧，删除这些弧，共删至多 $n$ 条；删除不会产生新圈。得到图 $H$，弧数 $m_H\ge(1-\delta)pN-n$，有向围长至少 $g$，且每个次序的前向弧数仍不超过原图相应值。因此
\[
 \mathrm{OPT}(H)\le \frac{m_H}{2}+\frac n2+\frac{\delta pN}{2}
 \le\left(\frac12+2\delta\right)m_H
\]
对充分大 $n$ 成立：$n/(pN)\to0$，而 $\delta/[2(1-\delta)]<2\delta$，留下吸收 $n$ 项的正余量。

给 $H$ 的每条弧赋 $x_e=1-1/g$。所有有向圈长度 $j\ge g$，故其和为 $j(1-1/g)\le j-1$，这个分数向量可行。于是
\[
 \frac{\mathrm{LP}(H)}{\mathrm{OPT}(H)}
 \ge\frac{1-1/g}{1/2+2\delta}.
\]
先使 $g$ 任意大、$\delta$ 任意小，再取相应充分大 $n$，右边趋于2。故2是全部有限实例比值的上确界；这里的构造是存在性证明，不声称每个有限图都达到2。

(d) 对待判点先查 $0\le x_e\le1$，不满足则返回被违反的边界行。通过后令 $y_e=1-x_e\ge0$；圈约束等价于
\[
 \sum_{e\in C}y_e\ge1\quad(\text{每个有向圈 }C).
\]
对每个起点运行非负权最短路，求全部距离 $d_y(v,u)$。对每条弧 $e=u\to v$ 计算 $y_e+d_y(v,u)$，不可达时记为无穷。所有这些值的最小值，就是最小 $y$ 权有向圈的权重：任一圈去掉一条弧后给一条从该弧终点回起点的路，因此候选最小值不大于最小圈值；反过来，若回路距离有限，取简单最短路（零权圈可删），加 $e$ 得一个简单圈，给同值候选。因此两边相等。

若最小值小于1，恢复相应圈 $C$，返回 $\sum_{e\in C}x_e\le |C|-1$，它包含整个可行多面体而严格排除当前点。若无圈或最小值至少1，全部圈约束均满足，正确确认可行。负容量/负 $y$ 的情况已由第一步排除，不能对未经查界的输入直接使用 Dijkstra。

邻接表、索引二叉堆的全源 Dijkstra 需 $O(n(n+m)\log n)$ 次基本运算，扫描候选和恢复圈再需 $O(n+m)$。输出行只有0、1系数与整数 $|C|-1$，位长为多项式；有理距离由最多 $n-1$ 条边相加，位长亦为多项式。此多面体包含 $[1/4-1/(8m),\,1/4+1/(8m)]^m$：各坐标在0、1之间，长度至少2的圈左边至多 $|C|(1/4+1/(8m))\le3|C|/8\le |C|-1$。故无自环且 $m\ge1$ 时它是满维有界非空有理多面体。按本节作业3第6题所明确调用的强分离到精确优化基础定理，该预言机可多项式求出本 LP 的最优有理向量及值；直接返回这些弧变量即可，不需要从分数向量恢复拓扑序。只有整数可行向量才直接表示所求无圈子图。
\end{solution}
